\FloatBarrier
\section{Discussion}\label{sec:discussion}

\subsection{Why a negative result is informative here}\label{sec:whyinformative}

A negative result is interpretable only relative to the rule, the sensitivity of the test and the location of the error. The chronology is qualified rather than prospective: the Edwards H1/H2 gates were frozen and archived before their respective score passes, while the cod 5\% band and some comparator declarations were formalised after its first pass; none of the disclosed completions changes the cod verdict. Cod operating characteristics are estimated by simulation under the specified data-generating processes. J-17 has a different, non-equivalent diagnostic: the realized-flux score of 7.55\,ft conditions on target-year drivers unavailable at the origin and is neither a power estimate nor forecast-skill evidence. Proposition~\ref{prop:increment} gives an exact identity, and the rolling J-17 application shows a $+76.0$\,ft$^2$ substitution term partly offset by a $-27.5$\,ft$^2$ residual cross-term before clipping; the one-step output cap reduces the observed M2--persistence gap to $+41.0$\,ft$^2$ (SI S23). On cod, the deficit is a misspecification term concentrated in the 1991--1995 transitions, which no member of the tested class can produce at the recorded catch (Proposition~\ref{prop:repeller}). On J-17, the low persistence of the flux component helps explain why carrying last year's recharge performs poorly, but the exact score gap is not attributable to substitution alone. Any forecast using partial-year recharge would need a separately defined origin and evaluation; it is not evaluated here. A future cod test of the low-biomass pattern would likewise need an external regime boundary, consistent catch timing, and the declared score rule applied to both series.

\subsection{Protocols, bands and the AR(1)}\label{sec:protocoldiff}

The cod and Edwards point rules disagree about M1: its 0.39\,ft one-year margin passes the Edwards H1/H2 gate but not the cod 5\% two-horizon gate, because 2.96\% lies inside the band and the five-year ordering reverses. The Edwards bootstrap interval covers zero in both implementations and MAE is a tie, so the point-rule retention is weak evidence of material gain; we nevertheless report it as the frozen H1/H2 rule dictates, not as a general advantage of output-only models. The class decisions are a separate issue. M2m passes Pass~1 H1/H2 and is excluded from the count of additional function classes by the broader frozen specification; its $-0.95$\,ft margin over persistence is the only one-year benchmark margin that separates from noise. In Pass~2, Renso, Rprecip and combo pass the point gate against persistence and M1, but the archived class assessment demotes them as M2m-plus-intercept variants; Rar fails the point gate. Thus the empty Pass~2 structure set is a class outcome, not a score-gate failure. Reporting the score and the class decision separately preserves both the numerical result and the protocol's declared structural interpretation.

\subsection{Limitations}\label{sec:limits}

The test covers two systems, three series, two model classes and two horizons; it is a demonstration and a replication, not a survey. On cod, each specification uses a single retrospective vintage, which favours the models (the 1983--2015 overlap between vintages differs by 24.7\,kt on average and by 131\,kt in one year); catch enters as a point series although both assessments treat it as bounded and the collapse years are under-reported; forecasts are point forecasts scored by RMSE and MAE rather than predictive distributions; Specification A has few origins; period bands were chosen after seeing the data; the archived runner's catch-timing offset moves catch-supplied one-step margins by 20--75\,kt without moving the verdict; and the operating-characteristic results cover only the declared in-class truths and two specific out-of-class truths. On J-17, recharge and pumpage are total-area estimates that mix the San Antonio and Uvalde pools, recharge is estimated by the stream-loss method, pumpage includes unreported use and is partly an institutional response to head, the map is not a karst model, the sample has four short fixed windows, and the nClimDiv precipitation source file is not redistributed (its provenance is recorded, and the scored panel is archived). The identity in Proposition~\ref{prop:increment} is exact, but the half-autocorrelation shortcut additionally requires stationarity, a mean-zero mechanism and a zero residual cross-term. The J-17 application therefore uses the full origin-specific increment and reports both empirical terms; the realized-flux diagnostic conditions on future inputs, and the output clip changes the scored error gap. Remark~\ref{rem:crossover} is a separate stationary benchmark identity; its approximation on J-17 does not establish stationarity, and it does not describe the cod series. Neither system used a prospective registration, and no such status is claimed.

\section{Conclusions}\label{sec:conclusions}

Across two unrelated systems, process-based forecast ladders were evaluated against persistence and the training mean under distinct, system-specific protocols. No surplus-production module was retained on either Northern cod assessment series; the principal failures coincide with the 1991--1995 collapse, which the tested model class cannot generate at the recorded catch. On J-17, M2's persisted-flux forecast has 14.70\,ft RMSE versus 13.23\,ft for persistence. The one-year M2m and Pass~2 climate variants pass their point gates, but the broader frozen specification classifies them as variants of the same affine-AR(1) function class; the separate Pass~1 M1 point-rule retention has a margin within uncertainty. No additional Edwards structure is retained by the class screen. The exact rolling-origin identity makes the J-17 score comparison more precise: driver substitution contributes $+76.03$\,ft$^2$, residual covariance contributes $-27.49$\,ft$^2$, and the raw MSE gap is $+48.54$\,ft$^2$; the single upper output clip reduces the archived scored gap to $+40.99$\,ft$^2$. A 7.55-ft realized-flux result is conditional on future inputs and is not forecast skill. The half-autocorrelation threshold is a special case of the identity, not a substitute for its full terms; the separate stationary baseline comparison crosses at $h=2$ on the J-17 origins and does not describe cod. These results distinguish process fidelity, information available at forecast origin, point-score evidence and structural-class evidence. They support no general claim about other systems: each inference is bounded by the series, origins and model classes evaluated here.
